%! Author = Omar Iskandarani
%! Date = 2026-07-24
%! Non-release speculative appendix v0.3

\documentclass[11pt]{article}
\usepackage[T1]{fontenc}
\usepackage[utf8]{inputenc}
\usepackage[a4paper,margin=1in]{geometry}
\usepackage{amsmath,amssymb,amsthm,bm}
\usepackage{mathtools}
\usepackage{booktabs,tabularx,array}
\usepackage{hyperref}
\hypersetup{colorlinks=true,linkcolor=blue,citecolor=blue,urlcolor=blue}

\newtheorem{definition}{Definition}[section]
\newtheorem{hypothesis}{Hypothesis}[section]
\newcommand{\Z}{\mathbb Z}

\title{SST Non-Release Speculative Research Appendix v0.3\\
\large Composite Links, Phase Products, Modulo Classes, and Closed Modulation Ans\"atze}
\author{Omar Iskandarani\\Independent Researcher, Groningen, The Netherlands}
\date{2026-07-31}

\begin{document}
\maketitle

\begin{quote}
\textbf{Non-release status.}
This document records exploratory constructions, geometric intuitions, and
unvalidated ans\"atze.  Its contents are not part of the SST Canon, are not part
of the release research-track companion, do not define SST predictions, and
may be modified or rejected without requiring a Canon version change.  No item
in this appendix may be used to tune a Canon calibration, optimizer, null test,
or stopping rule before formal promotion.
\end{quote}

\section{Epistemic ledger and promotion rule}

\begin{center}
\begin{tabularx}{0.96\textwidth}{@{}>{\bfseries}p{0.37\textwidth}X@{}}
\toprule
Item & Non-release status \\
\midrule
Periodic exponential modulation & \texttt{[ANSATZ]}; one closed modulation family among many. \\
Torus link $T(6,9)$ & \texttt{[RESEARCH CANDIDATE]}; exact three-component topology, no established SST selection law. \\
Three components times three phase labels & \texttt{[STRUCTURAL ANSATZ]}; naturally $\Z_3\times\Z_3$, not automatically $\Z_9$. \\
Order-nine carry dynamics & \texttt{[SPECULATIVE]}; requires an explicit physical coupling law. \\
Doubling map modulo nine & \texttt{[EXACT DISCRETE MATHEMATICS / PHYSICAL STATUS OPEN]}. \\
Rodin-style toroidal projection & \texttt{[INSPIRATION ONLY]}; geometry depends on representation choices. \\
Translation-induced trefoil flattening & \texttt{[RETIRED]}; not required by the action--phase mass--shell clock route. \\
Translation-induced stretching / uniform core thinning & \texttt{[RETIRED]}; no momentum--shape coupling has been derived. \\
Foliation-ordered Bell/Born programme & \texttt{[SPECULATIVE TEST PROGRAMME]}; no signalling, measurement independence, and anti-tuning gates mandatory. \\
\bottomrule
\end{tabularx}
\end{center}

Promotion into the release research track requires, at minimum:
\begin{enumerate}
    \item independently defined physical variables and an evolution law;
    \item dimensional consistency and conservation-law compatibility;
    \item comparison against alternative moduli, geometries, and phase laws;
    \item a preregistered falsifier and a minimal analytic or numerical test;
    \item no use of the candidate pattern in calibration or model selection;
    \item reproducible improvement over simpler control models.
\end{enumerate}

\section{Retired translation-induced shape hypotheses}

The exploratory suggestions that a translating electron trefoil must flatten,
stretch, or uniformly thin as its speed approaches \(c\) are retired.  The
v0.8.28 release route instead uses an internal action--phase Hamiltonian,
\begin{equation}
    H(P,I)=\sqrt{P^2c^2+E_0^2(I)},
\end{equation}
which produces the conditional phase-rate law
\begin{equation}
    \dot\theta=\Omega_0\,\frac{E_0}{H}=\frac{\Omega_0}{\gamma}
\end{equation}
without introducing a shape variable.  If shape coordinates \(q\) are added
only through \(E_0(I,q)\), then
\begin{equation}
    \partial_qH=\frac{E_0}{H}\,\partial_qE_0,
\end{equation}
so the rest-shape extrema are unchanged at finite momentum.  A future
translation-dependent deformation may be reconsidered only if an independently
derived coupling \(H_{Pq}\neq0\) is found.  The retired hypotheses may not be
used in calibration, particle-size estimates, or relativistic predictions.

\section{Periodic exponential modulation}

A bounded periodic radius modulation may be written as
\begin{equation}
    \rho(u)
    =
    \rho_0\exp\!\left[\eta\sin(nu)\right],
    \qquad
    n\in\Z,
    \label{eq:spec_periodic_exponential_radius}
\end{equation}
with a torus-knot or torus-link parameterization
\begin{equation}
    \mathbf X_{p,q,n,\eta}(u)
    =
    \begin{pmatrix}
    [R+\rho(u)\cos(qu)]\cos(pu)\\
    [R+\rho(u)\cos(qu)]\sin(pu)\\
    \rho(u)\sin(qu)
    \end{pmatrix},
    \qquad 0\le u<2\pi.
    \label{eq:spec_periodically_modulated_torus_embedding}
\end{equation}
For integer $p,q,n$, the parameterization is closed.  Closure alone does not
establish preservation of knot type, finite reach, adequate core clearance,
or dynamical stability.  The parameter $\eta$ must remain free in the primary
test.  Assignments such as $\eta=\ln\phi$ are secondary hypotheses and are not
permitted in an optimizer before the free-$\eta$ comparison.

The mandatory controls include an unmodulated embedding, ordinary sinusoidal
or Fourier modulation, a spline family with comparable degrees of freedom,
and equalized length or energy constraints.  The exponential family is
promoted only if it produces a reproducible dynamical advantage not explained
by parameter count or geometric clearance.

\section{The $T(6,9)$ three-trefoil candidate}

For the torus link $T(p,q)$, the number of components is
\begin{equation}
    d=\gcd(p,q).
\end{equation}
Consequently,
\begin{equation}
    \gcd(6,9)=3,
\end{equation}
and each component has reduced torus-knot type
\begin{equation}
    T\!\left(\frac{6}{3},\frac{9}{3}\right)
    =T(2,3),
\end{equation}
so $T(6,9)$ consists of three trefoil components.  With consistent orientation,
the pairwise linking number of parallel torus-link components is
\begin{equation}
    \operatorname{Lk}_{ij}
    =
    \frac{pq}{d^2}
    =6.
\end{equation}
These are topological facts.  They do not imply that $T(6,9)$ is selected by
Euler or Biot--Savart dynamics, that it is more stable than alternative
three-component links, or that it realizes a particle state.

\begin{hypothesis}[Composite-link candidate]
After a single resolved trefoil reaches at least the declared numerical
KAM--3 gate under a fixed protocol, $T(6,9)$ may be compared with control
three-component links under matched core radius, circulation, total length,
helicity metadata, and minimum clearance.
\end{hypothesis}

\section{Three-by-three phase labels}

Label the three components by $a\in\Z_3$ and three local phase sectors by
$k\in\Z_3$.  The nine labels form
\begin{equation}
    (a,k)\in\Z_3^{\rm component}\times\Z_3^{\rm phase},
    \qquad
    \left|\Z_3\times\Z_3\right|=9.
    \label{eq:spec_z3_product}
\end{equation}
This product group is not cyclic:
\begin{equation}
    \Z_3\times\Z_3\not\cong\Z_9,
\end{equation}
because every nonzero element of $\Z_3\times\Z_3$ has order three, whereas
$\Z_9$ contains elements of order nine.

A genuine order-nine orbit therefore requires an additional successor law.
One possible bookkeeping map is the base-three carry operation
\begin{equation}
    S(a,k)
    =
    \begin{cases}
    (a+1,k), & a=0,1,\\
    (0,k+1\!\!\pmod 3), & a=2,
    \end{cases}
    \qquad S^9=\operatorname{id}.
    \label{eq:spec_order_nine_carry}
\end{equation}
Equation~\eqref{eq:spec_order_nine_carry} is only a discrete indexing rule
until a physical interaction derives it.

\section{Doubling modulo nine}

The map
\begin{equation}
    x_{m+1}=2x_m\pmod 9
    \label{eq:spec_mod9_doubling}
\end{equation}
is an exact finite dynamical system.  Its orbit decomposition is
\begin{align}
    1&\rightarrow2\rightarrow4\rightarrow8\rightarrow7\rightarrow5\rightarrow1,
    \\
    3&\leftrightarrow6,
    \\
    0&\rightarrow0.
\end{align}
The six-cycle is the multiplicative unit group
\begin{equation}
    \Z_9^{\times}=\{1,2,4,5,7,8\},
    \qquad
    \operatorname{ord}_9(2)=6.
\end{equation}
This is rigorous number theory.  The familiar decimal digital-root
representation follows from $10\equiv1\pmod9$ and is therefore partly tied to
the choice of base ten.  A physical role for Eq.~\eqref{eq:spec_mod9_doubling}
requires an independently derived nine-state variable and a reason that one
physical update acts as multiplication by two modulo nine.

\section{Foliation-ordered Bell, Born, and no-signalling programme}
\label{sec:spec_bell_born_programme}

\textbf{[NON-RELEASE SPECULATIVE TEST PROGRAMME].}
A preferred ordering variable does not by itself reproduce quantum
nonlocality, the Born rule, or operational no-signalling.  This programme
records the minimum gates for testing such a construction without building the
quantum target into the kernel.  No result in this section may be used to
calibrate the release Canon.

\subsection*{Hidden-variable and setting-independence declaration}

Let \(x,y\) denote the two measurement settings, \(a,b\in\{-1,+1\}\) the
outcomes, and \(\lambda\) the complete model state.  The primary branch must
state whether it assumes measurement independence,
\begin{equation}
    P(\lambda\mid x,y)=P(\lambda).
    \label{eq:spec_measurement_independence}
\end{equation}
If Eq.~\eqref{eq:spec_measurement_independence} is abandoned, that branch must
be labelled explicitly as measurement-dependent or superdeterministic and must
quantify the required setting--state mutual information.  It may not be
presented as an ordinary local-hidden-variable resolution of Bell's theorem.

\subsection*{No-signalling gate}

Operational no-signalling requires the observable marginals
\begin{align}
    \sum_bP(a,b\mid x,y)&=P(a\mid x),
    \\
    \sum_aP(a,b\mid x,y)&=P(b\mid y),
    \label{eq:spec_no_signalling_marginals}
\end{align}
for every tested setting pair, not only after averaging over a target data set.
The preregistered residuals are
\begin{align}
    \mathcal R_{A\leftarrow B}
    &=
    \max_{a,x,y,y'}
    \left|
        P(a\mid x,y)-P(a\mid x,y')
    \right|,
    \\
    \mathcal R_{B\leftarrow A}
    &=
    \max_{b,y,x,x'}
    \left|
        P(b\mid x,y)-P(b\mid x',y)
    \right|.
    \label{eq:spec_no_signalling_residuals}
\end{align}
A foliation-ordered internal update is admissible only if these operational
residuals converge to zero within statistically declared uncertainty.

\subsection*{CHSH gate and anti-target-tuning rule}

Define
\begin{align}
    S
    =
    E_{00}+E_{01}+E_{10}-E_{11}.
    \label{eq:spec_chsh_parameter}
\end{align}
Local measurement-independent hidden-variable models satisfy \(|S|\le2\),
whereas quantum theory satisfies the Tsirelson bound \(|S|\le2\sqrt2\)
\cite{Bell1964,CHSH1969,Tsirelson1980}.  The model must reproduce the full
setting-dependent correlation curve on held-out settings, not merely tune one
four-setting combination to \(S>2\).  Mandatory controls include detector
inefficiency, time-window selection, postselection, setting predictability,
memory effects, and confidence intervals consistent with modern loophole-free
tests \cite{Hensen2015,Giustina2015,Shalm2015}.

\subsection*{Born-frequency gate}

For an independently defined state amplitude \(\psi_i\), the empirical
frequency \(f_i^{(N)}\) must approach
\begin{equation}
    p_i=|\psi_i|^2
    \label{eq:spec_born_target}
\end{equation}
without using \(|\psi_i|^2\) in the microscopic transition probabilities,
noise schedule, stopping rule, or sample weighting.  A primary diagnostic is
\begin{equation}
    D_{\rm KL}
    \!\left(
        f^{(N)}\middle\| |\psi|^2
    \right)
    \longrightarrow0
    \qquad(N\to\infty),
    \label{eq:spec_born_kl_gate}
\end{equation}
across held-out amplitudes, bases, initial conditions, and coarse-graining
choices.  Reproducing one two-slit histogram is insufficient.  The derivation
must also explain basis dependence and composition of independent systems
rather than assuming the Hilbert-space measure \cite{Born1926}.

\subsection*{Promotion blockers}

The programme remains rejected for release use if any of the following occurs:
\begin{enumerate}
    \item the nonlocal kernel contains the desired quantum correlator or Born
    weight as an input;
    \item signalling cancellation occurs only after parameter fine tuning or
    postselection;
    \item the CHSH violation disappears on held-out settings or after closure of
    detector/time-window loopholes;
    \item measurement dependence is present but unreported;
    \item the frequency measure changes under an arbitrary reparameterization
    or coarse-graining choice;
    \item the same parameters are used both to choose the model and to claim its
    validation.
\end{enumerate}

\paragraph{Child-level analogy.}
Two distant boxes may share a hidden instruction sheet, but they may not use a
secret telephone after the questions are chosen.  Getting one clever score is
not enough: the boxes must give the right pattern for many unseen questions,
and neither box may reveal which question the other box received.

\section{Separation from the KAM release programme}

The v0.8.25 KAM release programme tests relative equilibria, reduced Floquet
ellipticity, Krein signatures, nonlinear twist, finite-order non-resonance,
Poincar\'e maps, action drift, and invariant-torus survival.  The constructions
in this appendix may not be used to preselect a stable state.  The valid order
is
\begin{equation}
    \text{unbiased KAM search}
    \longrightarrow
    \text{observed structure}
    \longrightarrow
    \text{post-hoc comparison with speculative candidates}.
\end{equation}
A candidate is rejected when its apparent advantage disappears under matched
controls, resolution refinement, alternative parameterizations, or removal of
base-dependent encoding choices.

\paragraph{Child-level analogy.}
The appendix is a sketchbook, not the construction manual.  A beautiful gear
pattern can suggest what to test, but the machine must still work when the
pattern is hidden from the engineer.

\begin{thebibliography}{99}

\bibitem{Born1926}
M.~Born,
1926,
``Zur Quantenmechanik der Sto{\ss}vorg{\"a}nge,''
\emph{Zeitschrift f{\"u}r Physik}, vol.~37, pp.~863--867,
\href{https://doi.org/10.1007/BF01397477}{doi:10.1007/BF01397477}.

\bibitem{Bell1964}
J.~S. Bell,
1964,
``On the Einstein Podolsky Rosen paradox,''
\emph{Physics Physique Fizika}, vol.~1, pp.~195--200,
\href{https://doi.org/10.1103/PhysicsPhysiqueFizika.1.195}{doi:10.1103/PhysicsPhysiqueFizika.1.195}.

\bibitem{CHSH1969}
J.~F. Clauser, M.~A. Horne, A.~Shimony, and R.~A. Holt,
1969,
``Proposed experiment to test local hidden-variable theories,''
\emph{Physical Review Letters}, vol.~23, pp.~880--884,
\href{https://doi.org/10.1103/PhysRevLett.23.880}{doi:10.1103/PhysRevLett.23.880}.

\bibitem{Tsirelson1980}
B.~S. Tsirelson,
1980,
``Quantum generalizations of Bell's inequality,''
\emph{Letters in Mathematical Physics}, vol.~4, pp.~93--100,
\href{https://doi.org/10.1007/BF00417500}{doi:10.1007/BF00417500}.

\bibitem{Hensen2015}
B.~Hensen et al.,
2015,
``Loophole-free Bell inequality violation using electron spins separated by 1.3 kilometres,''
\emph{Nature}, vol.~526, pp.~682--686,
\href{https://doi.org/10.1038/nature15759}{doi:10.1038/nature15759}.

\bibitem{Giustina2015}
M.~Giustina et al.,
2015,
``Significant-loophole-free test of Bell's theorem with entangled photons,''
\emph{Physical Review Letters}, vol.~115, 250401,
\href{https://doi.org/10.1103/PhysRevLett.115.250401}{doi:10.1103/PhysRevLett.115.250401}.

\bibitem{Shalm2015}
L.~K. Shalm et al.,
2015,
``Strong loophole-free test of local realism,''
\emph{Physical Review Letters}, vol.~115, 250402,
\href{https://doi.org/10.1103/PhysRevLett.115.250402}{doi:10.1103/PhysRevLett.115.250402}.\bibitem{Rolfsen1976Knots}
D.~Rolfsen,
1976,
\emph{Knots and Links},
Publish or Perish,
reprinted by AMS Chelsea Publishing,
\url{https://bookstore.ams.org/chel-346}.

\bibitem{HardyWright2008NumberTheory}
G.~H.~Hardy and E.~M.~Wright,
2008,
\emph{An Introduction to the Theory of Numbers},
6th ed., Oxford University Press,
\url{https://global.oup.com/academic/product/an-introduction-to-the-theory-of-numbers-9780199219865}.

\bibitem{Moffatt1969Knottedness}
H.~K.~Moffatt,
1969,
``The Degree of Knottedness of Tangled Vortex Lines,''
\emph{Journal of Fluid Mechanics}, vol.~35, no.~1, pp.~117--129,
\href{https://doi.org/10.1017/S0022112069000991}{doi:10.1017/S0022112069000991}.

\end{thebibliography}
\end{document}
